\documentclass[10pt,a4paper,twocolumn]{article}
\usepackage{amsmath, amssymb, amsfonts}
\usepackage{mathtools}
\usepackage{geometry}
\usepackage[hidelinks]{hyperref}
\geometry{margin=0.75in}
\setlength{\columnsep}{0.25in}
\usepackage{parskip}
\usepackage{xcolor}
\usepackage{booktabs}
\usepackage{longtable}
\usepackage{caption}
\usepackage[raggedright]{titlesec}
\titlelabel{\thetitle.\hspace{0.5em}}
% abstract heading: centred, normal size
\renewenvironment{abstract}{\par\centerline{\bfseries Abstract}\par\vspace{0.3em}}{\par\vspace{0.5em}}
% full-width tables: fill pages from the top, allow more of them next to text
\makeatletter
\setlength{\@dblfptop}{0pt}
\setlength{\@dblfpsep}{14pt}
\setlength{\@dblfpbot}{0pt plus 1fil}
\makeatother
\setcounter{dbltopnumber}{3}
\renewcommand{\dbltopfraction}{0.85}
\renewcommand{\dblfloatpagefraction}{0.6}
\captionsetup[table]{font=small,labelfont=bf,skip=6pt,position=below}
\newsavebox{\captabbox}
\newlength{\captabwidth}
% table with its caption below, the caption as wide as the table (at least 0.6\linewidth)
\newenvironment{captab}[1]{\def\captabtext{#1}\begin{lrbox}{\captabbox}}{\end{lrbox}\setlength{\captabwidth}{\wd\captabbox}\ifdim\captabwidth<0.6\linewidth\setlength{\captabwidth}{0.6\linewidth}\fi\begin{minipage}{\captabwidth}\centering\usebox{\captabbox}\captionof{table}{\captabtext}\end{minipage}}

\newcommand{\R}{\mathbb{R}}

\begin{document}

	$p = \mathbb{P}(z = 1)$ is the fraction of behavior prompts. $\delta = \mu_1 - \mu_0$ is the difference of the class means, where $\mu_i = \mathbb{E}[x \mid z = i]$. $\mathrm{Cov}(x, z)$ is the cross-covariance, a vector in $\R^d$. $z \in \{0, 1\}$, where $z = 0$ is a neutral prompt and $z = 1$ is a behavior prompt, so $z \sim \mathrm{Bernoulli}(p)$.

	$$\begin{aligned}
		\mathrm{Cov}(x, z) & = \mathbb{E} \left[ (x - \mathbb{E}[x]) (z - \mathbb{E}[z]) \right] \\
		& = \mathbb{E}[x z] - \mathbb{E}[x] \mathbb{E}[z]
	\end{aligned}$$

	$$\mathbb{E}[z] = 1 \cdot p + 0 \cdot (1 - p) = p$$

	$$\begin{aligned}
		\mathbb{E}[x] & = \mathbb{P}(z = 1) \mathbb{E}[x \mid z = 1] + \mathbb{P}(z = 0) \mathbb{E}[x \mid z = 0] \\
		& = p \mathbb{E}[x \mid z = 1] + (1 - p) \mathbb{E}[x \mid z = 0] \\
		& = p \mu_1 + (1 - p) \mu_0 = \mu
	\end{aligned}$$

	$$\begin{aligned}
		\mathbb{E}[x z] & = p \mathbb{E}[x z \mid z = 1] + (1 - p) \mathbb{E}[x z \mid z = 0] \\
		& = p \mu_1
	\end{aligned}$$
	since the second term is $0$.

	$$\mathrm{Cov}(x, z) = \mathbb{E}[x z] - \mathbb{E}[x] \mathbb{E}[z] = p \mu_1 - \mu p = p (\mu_1 - \mu)$$

	Since $\mu = p \mu_1 + (1 - p) \mu_0$,
	$$\begin{aligned}
		p (\mu_1 - \mu) & = p \left( \mu_1 - p \mu_1 - (1 - p) \mu_0 \right) \\
		& = p (1 - p) (\mu_1 - \mu_0) = p (1 - p) \delta
	\end{aligned}$$
	so
	$$\mathrm{Cov}(x, z) = p (1 - p) \delta$$

	By the law of total covariance,
	$$\begin{aligned}
		\mathrm{Cov}(x) = {} & p \mathrm{Cov}(x \mid z = 1) + (1 - p) \mathrm{Cov}(x \mid z = 0) \\
		& + \mathrm{Cov}(\mathbb{E}[x \mid z])
	\end{aligned}$$

	$$\mathbb{E}_z \left[ \mathbb{E}[x \mid z] \right] = (1 - p) \mu_0 + p \mu_1 = \mu$$

	$$\begin{aligned}
		\mu_0 - \mu & = p (\mu_0 - \mu_1) \\
		\mu_1 - \mu & = (1 - p) (\mu_1 - \mu_0)
	\end{aligned}$$

	$$\begin{aligned}
		\mathrm{Cov}(\mathbb{E}[x \mid z]) = {} & \mathbb{P}(z = 0) (\mu_0 - \mu) (\mu_0 - \mu)^\top \\
		& + \mathbb{P}(z = 1) (\mu_1 - \mu) (\mu_1 - \mu)^\top \\
		= {} & (1 - p) p^2 (\mu_0 - \mu_1) (\mu_0 - \mu_1)^\top \\
		& + p (1 - p)^2 (\mu_1 - \mu_0) (\mu_1 - \mu_0)^\top \\
		= {} & \left( (1 - p) p^2 + p (1 - p)^2 \right) \delta \delta^\top \\
		= {} & p (1 - p) \left( p + (1 - p) \right) \delta \delta^\top \\
		= {} & p (1 - p) \delta \delta^\top
	\end{aligned}$$

	So
	$$\mathrm{Cov}(x) = \mathbb{E}_z \left[ \mathrm{Cov}(x \mid z) \right] + p (1 - p) \delta \delta^\top$$
	where
	$$\begin{aligned}
		\mathbb{E}_z \left[ \mathrm{Cov}(x \mid z) \right] = {} & p \mathrm{Cov}(x \mid z = 1) \\
		& + (1 - p) \mathrm{Cov}(x \mid z = 0)
	\end{aligned}$$

	If $\mathrm{Cov}(x) \in \mathrm{GL}_d(\R)$, its range is all of $\R^d$, but it can be singular (for example, with fewer samples than dimensions). If $\mathrm{Cov}(x) v = 0$ for some $v$, then $\delta^\top v = 0$, so
	$$\delta \in \mathrm{range}(\mathrm{Cov}(x)) = \{ \mathrm{Cov}(x) u : u \in \R^d \}$$
	A system $A u = y$ is solvable if and only if $y \in \mathrm{range}(A)$, so I need to check whether $\delta \in \mathrm{range}(\mathrm{Cov}(x))$.

	Well, we know that $\mathrm{Cov}(x)$ is symmetric, so its range is the orthogonal complement of its kernel:
	$$\mathrm{range}(\mathrm{Cov}(x)) = \ker(\mathrm{Cov}(x))^\perp$$
	so we need to show that $v^\top \delta = 0$ for all $v \in \ker(\mathrm{Cov}(x))$.

	Assume $\mathrm{Cov}(x) v = 0$. Then
	$$\begin{aligned}
		0 & = v^\top \mathrm{Cov}(x) v \\
		& = v^\top \mathbb{E}_z \left[ \mathrm{Cov}(x \mid z) \right] v + p (1 - p) (v^\top \delta)^2
	\end{aligned}$$
	$\mathbb{E}_z \left[ \mathrm{Cov}(x \mid z) \right]$ is a weighted sum of two covariance matrices with non-negative weights, so it is positive semi-definite, and $v^\top \mathbb{E}_z \left[ \mathrm{Cov}(x \mid z) \right] v \ge 0$. Since $p (1 - p) > 0$, we have $p (1 - p) (v^\top \delta)^2 \ge 0$. Both terms are non-negative, so for them to sum up to $0$, both must be $0$. Hence $v^\top \delta = 0$, and $\delta \in \mathrm{range}(\mathrm{Cov}(x))$.

	Inspired by LEACE, I will try to prove the following result for $r(x) = P x + b$:
	$$\begin{aligned}
		& \mathrm{Cov}(A r(x), z) = 0 \\
		& \iff A P \delta = 0 \\
		& \iff \mathbb{E}[A r(x) \mid z = 1] = \mathbb{E}[A r(x) \mid z = 0]
	\end{aligned}$$
	We know that $\mathbb{E}[z] = p$, $\mathbb{E}[x] = p \mu_1 + (1 - p) \mu_0 = \mu$ and $\mathbb{E}[x z] = p \mu_1$. Recall:
	$$\mathrm{Cov}(u, y) = \mathbb{E}[u y] - \mathbb{E}[u] \mathbb{E}[y]$$

	Let $u := A r(x) = A P x + A b$. Then
	$$\begin{aligned}
		\mathbb{E}[u z] & = \mathbb{E}[A P x z + A b z] = A P \mathbb{E}[x z] + A b \mathbb{E}[z] \\
		& = p A P \mu_1 + p A b
	\end{aligned}$$
	$$\mathbb{E}[u] = A P \mathbb{E}[x] + A b = A P \mu + A b$$
	$$\mathbb{E}[u] \mathbb{E}[z] = p A P \mu + p A b$$
	So
	$$\begin{aligned}
		\mathrm{Cov}(A r(x), z) & = p A P \mu_1 + p A b - (p A P \mu + p A b) \\
		& = p A P (\mu_1 - \mu)
	\end{aligned}$$
	Since $\mu_1 - \mu = (1 - p) \delta$,
	$$\mathrm{Cov}(A r(x), z) = p (1 - p) A P \delta$$
	Also, $p (1 - p) A P \delta = A P \mathrm{Cov}(x, z)$, since $\mathrm{Cov}(x, z) = p (1 - p) \delta$. Since $p (1 - p) > 0$,
	$$\mathrm{Cov}(A r(x), z) = 0 \iff A P \delta = 0$$

	By linearity, $\mathbb{E}[A r(x) \mid z = i] = A P \mu_i + A b$, so
	$$\begin{aligned}
		\mathbb{E}[A r(x) \mid z = 1] - \mathbb{E}[A r(x) \mid z = 0] & = A P (\mu_1 - \mu_0) \\
		& = A P \delta
	\end{aligned}$$
	and therefore
	$$A P \delta = 0 \iff \mathbb{E}[A r(x) \mid z = 1] = \mathbb{E}[A r(x) \mid z = 0]$$

	Okay, $A P \delta = 0$ restricts $P$ and says nothing about $b$, so right now many edits are feasible, and the cost $J$ (how far the edit moves the activations) is what can pick the best one.

	Let $r(x) = P x + b$ and $D = P - I$ (so basically, the distance to doing nothing), and let $M$ be symmetric positive definite (for the definition of the norm). The cost of the edit is
	$$J(P, b) = \mathbb{E} \left[ \| r(x) - x \|_M^2 \right]$$
	with the $M$-norm defined as $\| v \|_M^2 = v^\top M v$. I am searching for the closed form of $J(P, b)$.

	$$r(x) - x = D x + b = D (x - \mu) + (D \mu + b)$$
	where $D (x - \mu)$ is random with mean $0$ and $D \mu + b$ is constant.
	$$\begin{aligned}
		& \| D (x - \mu) + (D \mu + b) \|_M^2 \\
		& = \left( D (x - \mu) + (D \mu + b) \right)^\top M \left( D (x - \mu) + (D \mu + b) \right)
	\end{aligned}$$
	Matrix multiplication distributes over sums, so this equals
	$$\begin{aligned}
		& (D (x - \mu))^\top M (D (x - \mu)) + (D (x - \mu))^\top M (D \mu + b) \\
		& + (D \mu + b)^\top M (D (x - \mu)) + (D \mu + b)^\top M (D \mu + b)
	\end{aligned}$$
	Since $(X Y Z)^\top = Z^\top Y^\top X^\top$ and $M^\top = M$, and a scalar equals its own transpose,
	$$\begin{aligned}
		& (D (x - \mu))^\top M (D \mu + b) \\
		& = \left( (D (x - \mu))^\top M (D \mu + b) \right)^\top \\
		& = (D \mu + b)^\top M^\top D (x - \mu) \\
		& = (D \mu + b)^\top M D (x - \mu)
	\end{aligned}$$
	so the two middle terms sum to $2 (D \mu + b)^\top M D (x - \mu)$. Since $(D (x - \mu))^\top = (x - \mu)^\top D^\top$ and $(D \mu + b)^\top M (D \mu + b) = \| D \mu + b \|_M^2$,
	$$\begin{aligned}
		& \| D (x - \mu) + (D \mu + b) \|_M^2 \\
		& = (x - \mu)^\top D^\top M D (x - \mu) \\
		& + 2 (D \mu + b)^\top M D (x - \mu) + \| D \mu + b \|_M^2
	\end{aligned}$$

	Expectation is linear, $D$, $M$, $\mu$ and $b$ are constant, and $\mathbb{E}[x - \mu] = \mathbb{E}[x] - \mu = \mu - \mu = 0$, so
	$$\begin{aligned}
		J(P, b) = {} & \mathbb{E} \left[ (x - \mu)^\top D^\top M D (x - \mu) \right] \\
		& + 2 (D \mu + b)^\top M D \mathbb{E}[x - \mu] + \| D \mu + b \|_M^2 \\
		= {} & \mathbb{E} \left[ (x - \mu)^\top D^\top M D (x - \mu) \right] + \| D \mu + b \|_M^2
	\end{aligned}$$

	Now, the term inside the expectation looks like $\| v \|_{D^\top M D}^2$ with $v = x - \mu$. Recall that for vectors $a, c \in \R^d$ and a matrix $N \in \R^{d \times d}$ (not necessarily symmetric positive definite), $\mathrm{tr}(a c^\top) = c^\top a$, so
	$$\mathrm{tr}(N v v^\top) = \mathrm{tr}((N v) v^\top) = v^\top N v$$
	Since the trace is linear, it can be swapped with the expectation, and since $\mathbb{E}[(x - \mu)(x - \mu)^\top] = \mathrm{Cov}(x)$,
	$$\begin{aligned}
		& \mathbb{E} \left[ (x - \mu)^\top D^\top M D (x - \mu) \right] \\
		& = \mathbb{E} \left[ \mathrm{tr} \left( D^\top M D (x - \mu)(x - \mu)^\top \right) \right] \\
		& = \mathrm{tr} \left( D^\top M D \mathbb{E} \left[ (x - \mu)(x - \mu)^\top \right] \right) \\
		& = \mathrm{tr} \left( D^\top M D \mathrm{Cov}(x) \right)
	\end{aligned}$$
	So
	$$J(P, b) = \mathrm{tr} \left( D^\top M D \mathrm{Cov}(x) \right) + \| D \mu + b \|_M^2$$

	The first term does not depend on $b$. Since $\| D \mu + b \|_M^2 \ge 0$, it is minimized when it is equal to $0$, and because $M$ is positive definite, this happens exactly when $D \mu + b = 0$. Equivalently, the optimal offset is
	$$b^\star = - D \mu = (I - P) \mu$$
	so every optimal edit leaves the mean in place: $r(\mu) = P \mu + (I - P) \mu = \mu$.

	\section{Nullwalk is LEACE when $A = I$}

	Currently, our problem looks like this ($P \in \R^{d \times d}$, $b \in \R^d$):
	$$\min_{P, b} J(P, b) \text{ with constraint } \mathrm{Cov}(A r(x), z) = 0$$
	with $r(x) = P x + b$, $D = P - I$ and
	$$J(P, b) = \mathrm{tr} \left( D^\top M D \mathrm{Cov}(x) \right) + \| D \mu + b \|_M^2$$
	When $A = I$, we have $A r(x) = r(x)$, and the problem becomes
	$$\begin{aligned}
		& \min_{P, b} \mathbb{E} \left[ \| P x + b - x \|_M^2 \right] \\
		& \text{with constraint } \mathrm{Cov}(P x + b, z) = 0
	\end{aligned}$$
	Since $\mathrm{Cov}(A r(x), z) = 0 \iff A P \delta = 0$, with $A = I$ the constraint is equivalent to $P \delta = 0$.

	LEACE finds the optimal eraser in two steps: centered data first (Theorem 4.2), then the extension to uncentered data with an offset (Theorem 4.3). Our $x$ is not centered and our edit has an offset $b$, so Theorem 4.3 is the one we need.

	Theorem 4.3 of the LEACE paper in our notation: let $x$ and $Z$ be random vectors with values in $\R^d$ and $\R^K$, each with finite second moment. Let $M \in \R^{d \times d}$ be positive semi-definite, let $W = \left( \mathrm{Cov}(x)^{1/2} \right)^\dagger$, and let
	$$P_{W \mathrm{Cov}(x, Z)} = \left( W \mathrm{Cov}(x, Z) \right) \left( W \mathrm{Cov}(x, Z) \right)^\dagger$$
	be the orthogonal projection onto the column space of $W \mathrm{Cov}(x, Z)$. Define
	$$\begin{aligned}
		P^\star & = I - W^\dagger P_{W \mathrm{Cov}(x, Z)} W \\
		b^\star & = \mu - P^\star \mu
	\end{aligned}$$
	Then $(P^\star, b^\star)$ minimizes $\mathbb{E} \left[ \| P x + b - x \|_M^2 \right]$ subject to $\mathrm{Cov}(P x + b, Z) = 0$.

	For us $Z = z$, so $K = 1$, so Theorem 4.3 minimizes
	$$\mathbb{E} \left[ \| P x + b - x \|_M^2 \right] = J(P, b)$$
	and
	$$\mathrm{Cov}(P x + b, Z) = 0 \text{ reads } \mathrm{Cov}(P x + b, z) = 0$$

	Two optimization problems with the same objective and the same feasible set have the same minimizers, and by the theorem above, $(P^\star, b^\star)$ is a minimizer, so it also solves Nullwalk with $A = I$, and the edit it defines is
	$$\begin{aligned}
		r_{\mathrm{LEACE}}(x) & = P^\star x + b^\star \\
		& = x - W^\dagger P_{W \mathrm{Cov}(x, z)} W (x - \mu)
	\end{aligned}$$
	Consistency check for my optimal $b^\star$:
	$$\mu - P^\star \mu = - (P^\star - I) \mu = - D^\star \mu = b^\star$$

	For the justification of calling the constraint ``linear erasure'', encode $Z = (1 - z, z)^\top \in \{0, 1\}^2$ (each class with nonzero probability, since $p \in (0, 1)$). We use Definition 2.3 from the LEACE paper: $r(x)$ linearly guards $Z$ if, for every non-negative loss that is convex in its first argument, no linear predictor $\eta(r(x)) = c + V r(x)$, with $V \in \R^{K \times d}$, achieves a lower expected loss on $(r(x), Z)$ than the best constant predictor can.

	Also, by Theorem 3.4 of the LEACE paper: if each class probability is nonzero, then $\mathbb{E}[x \mid z = j] = \mathbb{E}[x]$ for every class $j$ if and only if $\mathrm{Cov}(x, Z) = 0$. And by Theorem 4.1 of the LEACE paper: for an affine $r(x) = P x + b$, $r(x)$ linearly guards $Z$ if and only if the columns of $\mathrm{Cov}(x, Z) \in \R^{d \times K}$ lie in $\ker(P)$.

	$\mathrm{Cov}(x, 1 - z) = - \mathrm{Cov}(x, z)$, so
	$$\mathrm{Cov}(x, Z) = \begin{pmatrix} - p (1 - p) \delta & p (1 - p) \delta \end{pmatrix}$$
	and both columns are nonzero multiples of $\delta$, so by Theorem 4.1 of the LEACE paper:
	$$\begin{aligned}
		r(x) \text{ linearly guards } Z & \iff P \delta = 0 \\
		& \iff \mathrm{Cov}(r(x), z) = 0
	\end{aligned}$$
	So Nullwalk with $A = I$ says: after the edit, no linear classifier can tell the two classes apart better than a constant predictor.

	Okay, now ``back to $A$''!

	I showed that $\mathrm{Cov}(A r(x), z) = p (1 - p) A P \delta$, and since $p (1 - p) > 0$, the Nullwalk constraint is equivalent to $A P \delta = 0$. The cost $J(P, b)$ does not contain $A$ at all, so $A$ enters the problem only through the constraint. I would like to know how much of $A$ the problem ``sees'', so I will rewrite the constraint with the kernel:
	$$\mathrm{Cov}(A r(x), z) = 0 \iff P \delta \in \ker(A)$$
	as follows trivially from $A P \delta = 0$: $A (P \delta) = 0$ and $\ker(A) = \{ v : A v = 0 \}$, hence $P \delta \in \ker(A)$.

	So $P \delta$ is the class difference that is left after the edit, and LEACE demands $P \delta = 0$, so nothing may be left. Nullwalk allows anything left over, as long as $A$ cannot see it.

	But do two different matrices with the same kernel define the same problem? Let $A \in \R^{k \times d}$ and $A' \in \R^{k' \times d}$ ($k$ and $k'$ may differ, but $A \ne A'$ anyway) with $\ker(A) = \ker(A')$. Denote the feasible set
	$$\begin{aligned}
		\mathcal{F}(A) = \{ & (P, b) \in \R^{d \times d} \times \R^d : \\
		& \mathrm{Cov}(A (P x + b), z) = 0 \}
	\end{aligned}$$
	so the Nullwalk problem for $A$ is
	$$\min_{(P, b) \in \mathcal{F}(A)} J(P, b)$$
	Since $\mathrm{Cov}(A (P x + b), z) = 0 \iff P \delta \in \ker(A)$, an equivalent way of writing $\mathcal{F}(A)$ is
	$$\mathcal{F}(A) = \{ (P, b) : P \delta \in \ker(A) \}$$
	(every $b \in \R^d$ is allowed). Since $\ker(A) = \ker(A')$, $\mathcal{F}(A) = \mathcal{F}(A')$.

	Since $J(P, b)$ does not depend on the readout $A$,
	$$\min_{(P, b) \in \mathcal{F}(A)} J(P, b) = \min_{(P, b) \in \mathcal{F}(A')} J(P, b)$$
	so they minimize the same function over the same set, so they have the same minimizers and the same optimal value.

	If $\ker(A) = \{0\}$,
	$$\begin{aligned}
		\mathcal{F}(A) & = \{ (P, b) : P \delta \in \ker(A) \} \\
		& = \{ (P, b) : P \delta = 0 \}
	\end{aligned}$$
	which is identical to the $A = I$ case, so the solution is $r_{\mathrm{LEACE}}$.

	For $A$ with full row rank $k \le d$,
	$$\dim(\ker(A)) = d - \mathrm{rank}(A) = d - k$$
	so $\dim(\ker(A)) = 0 \implies k = d$, and hence $\ker(A) = \{0\} \iff k = d$.

	To summarize the results above (readout matrices with the same kernel define the same problem):
	\begin{enumerate}
		\item We can rescale, reorder or orthonormalize the rows of $A$ as we see fit, as these actions do not change the kernel of the matrix.
		\item Linearly dependent rows of $A$ change nothing.
	\end{enumerate}
	And summarizing $\ker(A) = \{0\} \iff k = d$:
	\begin{enumerate}
		\item If the readout has as many independent directions as the space itself, nothing is hidden from it, and Nullwalk gives the same as LEACE.
		\item This is not a direct corollary, but I wanted to use the J-lens as $A$, because I really enjoy the idea of the J-lens, but the full J-lens matrix is probably invertible, so it would be useless here (maybe the rows of $A$ can be the top-$k$ gradient directions, or the vectors of the refusal tokens?).
	\end{enumerate}

	This is also not a direct corollary of either, but when I was summarizing the results, I was thinking whether more readers mean a larger edit, so I will try to figure that out now.

	Let $\ker(A) \subseteq \ker(A')$, so $A'$ sees fewer directions than $A$.
	$$\begin{aligned}
		(P, b) \in \mathcal{F}(A) & \implies P \delta \in \ker(A) \subseteq \ker(A') \\
		& \implies (P, b) \in \mathcal{F}(A')
	\end{aligned}$$
	so $\mathcal{F}(A) \subseteq \mathcal{F}(A')$. Minimizing $J$ over the larger set gives
	$$J^\star_{A'} = \inf_{\mathcal{F}(A')} J \le \inf_{\mathcal{F}(A)} J = J^\star_A$$
	Corollary: $\ker(I) = \{0\} \subseteq \ker(A')$ for every $A'$, hence
	$$J^\star_{A'} \le J^\star_{\mathrm{LEACE}}$$
	Because I used $\inf$, there is no assumption that the minimizer exists, and for $A = I$ a minimizer does exist by Theorem 4.3 of the LEACE paper.

	So adding rows to $A$ shrinks $\ker(A)$ (if they are linearly independent from the existing rows), so the optimal edit can only grow. More rows guard against more readers, but bring the edit closer to LEACE.

	Okay, now I want the LEACE formula adjusted for a binary label. From the LEACE paper: $W = \left( \mathrm{Cov}(x)^{1/2} \right)^\dagger$ and
	$$r_{\mathrm{LEACE}}(x) = x - W^\dagger P_{W \mathrm{Cov}(x, z)} W (x - \mu)$$

	I did encounter pseudoinverses during my undergraduate studies, but I do not have much experience with them. Every real symmetric matrix is orthogonally diagonalizable, and since $\mathrm{Cov}(x)$ is real symmetric (and positive semi-definite, but that is irrelevant here),
	$$\mathrm{Cov}(x) = Q \Lambda Q^\top$$
	with $\Lambda = \mathrm{diag}(\lambda_1, \dots, \lambda_d)$, $\lambda_i \ge 0$ and $Q^\top Q = I$.

	Pseudoinverses of a diagonal matrix are easy: $\Lambda^\dagger$ replaces each nonzero $\lambda_i$ by $1 / \lambda_i$ (and leaves $0$ as $0$), so it undoes scaling. Conjugating by an orthogonal $Q$ does not change this: $(Q B Q^\top)^\dagger = Q B^\dagger Q^\top$.

	Whitening from LEACE:
	$$W = Q (\Lambda^{1/2})^\dagger Q^\top$$
	(divide each coordinate by its standard deviation, so every direction ends up with variance $1$). Whitening twice is the pseudoinverse of the covariance: $\left( 1 / \sqrt{\lambda_i} \right)^2 = 1 / \lambda_i$, so
	$$\begin{aligned}
		W W & = Q (\Lambda^{1/2})^\dagger (\Lambda^{1/2})^\dagger Q^\top \\
		& = Q \Lambda^\dagger Q^\top = \mathrm{Cov}(x)^\dagger
	\end{aligned}$$
	and
	$$W^\top W = W W$$

	The pseudoinverse of a pseudoinverse is the original matrix, so
	$$W^\dagger = \mathrm{Cov}(x)^{1/2} = Q \Lambda^{1/2} Q^\top$$
	$$W^\dagger W = Q \mathrm{diag}(\mathbf{1}_{\lambda_i > 0}) Q^\top =: \Pi$$
	Each $\sqrt{\lambda_i}$ cancels against $1 / \sqrt{\lambda_i}$, so $\Pi$ keeps the directions with positive variance and zeroes the rest. It is the orthogonal projection onto $\mathrm{range}(\mathrm{Cov}(x))$.

	Applying $W$ and then $W^\dagger$ therefore loses only the part of a vector that lies in zero-variance directions. We proved that $\delta \in \mathrm{range}(\mathrm{Cov}(x))$, so $\delta$ has no such part:
	$$W^\dagger W \delta = \Pi \delta = \delta$$

	For a binary label, $\mathrm{Cov}(x, z) = p (1 - p) \delta$ is one column. Let $u := W \delta$, so $u$ is the mean difference measured in whitened coordinates, and
	$$W \mathrm{Cov}(x, z) = p (1 - p) u$$
	Since $W$ kills only zero-variance directions ($\ker(W) = \ker(\mathrm{Cov}(x))$), and we know that $\delta \in \mathrm{range}(\mathrm{Cov}(x)) = \ker(\mathrm{Cov}(x))^\perp$ and $\delta \ne 0$, we get $\delta \notin \ker(W)$, so $u \ne 0$.

	Following LEACE, we project onto that one direction:
	$$P_{W \mathrm{Cov}(x, z)} = (p (1 - p) u) (p (1 - p) u)^\dagger$$
	is the orthogonal projection onto the column space of $p (1 - p) u$, which is the line through $u$. For $v \ne 0$ and $v \in \R^d$,
	$$v^\dagger = \frac{v^\top}{v^\top v}$$
	since this row vector satisfies $v^\dagger v = 1$ (and the four Moore--Penrose conditions). So
	$$\begin{aligned}
		P_{W \mathrm{Cov}(x, z)} & = p (1 - p) u \frac{p (1 - p) u^\top}{(p (1 - p))^2 u^\top u} \\
		& = \frac{u u^\top}{u^\top u}
	\end{aligned}$$

	Plugging in and un-whitening:
	$$W^\dagger P_{W \mathrm{Cov}(x, z)} W = \frac{(W^\dagger u)(u^\top W)}{u^\top u}$$
	$$W^\dagger u = W^\dagger W \delta = \delta$$
	$$u^\top W = \delta^\top W^\top W = \delta^\top \mathrm{Cov}(x)^\dagger$$
	Let $w := \mathrm{Cov}(x)^\dagger \delta$. $\mathrm{Cov}(x)^\dagger$ is symmetric (since $\mathrm{Cov}(x) = \mathrm{Cov}(x)^\top$), so $u^\top W = w^\top$, and
	$$u^\top u = \delta^\top W^\top W \delta = w^\top \delta$$
	$u \ne 0$, so this is positive. So
	$$W^\dagger P_{W \mathrm{Cov}(x, z)} W = \frac{\delta w^\top}{w^\top \delta}$$
	and
	$$r_{\mathrm{LEACE}}(x) = x - \frac{\delta w^\top}{w^\top \delta} (x - \mu)$$

	Every activation is moved along $\delta$ only: $x$ changes only by a multiple of $\delta$ (since $\frac{w^\top (x - \mu)}{w^\top \delta}$ is a scalar).

	For more intuition, rewrite
	$$r_{\mathrm{LEACE}}(x) = x - \delta \frac{w^\top (x - \mu)}{w^\top \delta}$$
	$w^\top (x - \mu)$ is a linear score of how ``behavior-like'' $x$ is. $w^\top \delta$ is the score difference between the two class means. Dividing by it calibrates the move so that the class difference is removed exactly.

	Recall that for any edit $r(x) = P x + b$ with $D = P - I$,
	$$J(P, b) = \mathrm{tr} \left( D^\top M D \mathrm{Cov}(x) \right) + \| D \mu + b \|_M^2$$
	$$\begin{aligned}
		r_{\mathrm{LEACE}}(x) & = x - \delta \frac{w^\top (x - \mu)}{w^\top \delta} \\
		& = \underbrace{\left( I - \frac{\delta w^\top}{w^\top \delta} \right)}_{P^\star} x + \underbrace{\frac{\delta w^\top}{w^\top \delta} \mu}_{b^\star}
	\end{aligned}$$
	so
	$$\begin{aligned}
		D^\star & = P^\star - I = - \frac{\delta w^\top}{w^\top \delta} \\
		b^\star & = - D^\star \mu
	\end{aligned}$$
	and $\| D^\star \mu + b^\star \|_M^2 = 0$.
	$$\begin{aligned}
		D^{\star \top} M D^\star & = \frac{w \delta^\top M \delta w^\top}{(w^\top \delta)^2} \\
		& = \frac{\| \delta \|_M^2}{(w^\top \delta)^2} w w^\top
	\end{aligned}$$
	where $\frac{\| \delta \|_M^2}{(w^\top \delta)^2}$ is a scalar.
	$$\mathrm{tr} \left( w w^\top \mathrm{Cov}(x) \right) = w^\top \mathrm{Cov}(x) w$$
	(I think this is called the cyclic property of the trace).
	$$\begin{aligned}
		J(P^\star, b^\star) & = \frac{\| \delta \|_M^2}{(w^\top \delta)^2} \mathrm{tr} \left( w w^\top \mathrm{Cov}(x) \right) \\
		& = \frac{\| \delta \|_M^2}{(w^\top \delta)^2} w^\top \mathrm{Cov}(x) w
	\end{aligned}$$
	Since $w = \mathrm{Cov}(x)^\dagger \delta$, $\mathrm{Cov}(x)^\dagger$ is symmetric, and $B^\dagger B B^\dagger = B^\dagger$,
	$$\begin{aligned}
		w^\top \mathrm{Cov}(x) w & = \delta^\top \mathrm{Cov}(x)^\dagger \mathrm{Cov}(x) \mathrm{Cov}(x)^\dagger \delta \\
		& = \delta^\top \mathrm{Cov}(x)^\dagger \delta = w^\top \delta
	\end{aligned}$$
	One factor $w^\top \delta$ cancels. $w^\top \delta = \| W \delta \|^2 > 0$, and we get
	$$J^\star_{\mathrm{LEACE}} = \frac{\| \delta \|_M^2}{w^\top \delta} = \frac{\| \delta \|_M^2}{\delta^\top \mathrm{Cov}(x)^\dagger \delta}$$

\end{document}
